\chapter{Hensel-lifting orthogonal idempotents (Stacks 0DXB fragment)}
\label{chap:Mathlib_RingTheory_Etale_HenselianIdempotentLift}

This chapter tracks the Mathlib-PR-target helper
\lean{Algebra.Etale.exists_completeOrthogonalIdempotents_lift_of_henselian}
extracted (during iter-048) from
\lean{exists_idempotent_lift_isolating_at_maximal} in
\texttt{Proetale/Mathlib/RingTheory/Etale/StrictlyHenselian.lean}.

\begin{theorem}
  \label{thm:etale-henselian-cop-lift}
  \lean{Algebra.Etale.exists_completeOrthogonalIdempotents_lift_of_henselian}
  \leanok
  Let $A$ be a henselian local ring with maximal ideal $\mathfrak{m}$
  and residue field $k = A/\mathfrak{m}$, and let $B$ be an étale
  $A$-algebra. Let $k \otimes_A B \simeq \prod_{i \in I} k_i$ be the
  residue product decomposition (each $k_i$ a finite separable
  extension of $k$, $I$ finite). Then the canonical orthogonal
  idempotents $\{e_i\}_{i \in I}$ of the product $\prod_{i \in I} k_i$
  lift to a complete orthogonal idempotent system
  $\{\tilde e_i\}_{i \in I}$ in $B$, in the sense that the image of
  $1 \otimes \tilde e_i$ in $k \otimes_A B$ corresponds to $e_i$
  under the product decomposition.
\end{theorem}

\medskip\noindent\textbf{Prerequisites.} The Lean signature for
\texttt{Algebra.Etale.exists\_completeOrthogonalIdempotents\_lift\_of\_henselian}
carries two additional instance binders beyond
\texttt{[HenselianLocalRing A]} and \texttt{[Algebra.Etale A B]},
namely
\[
  \texttt{[IsNoetherianRing A]} \qquad \text{and} \qquad
  \texttt{[Module.Finite A B]}.
\]
These hypotheses are \emph{not} consequences of
\texttt{HenselianLocalRing A} + \texttt{Algebra.Etale A B} alone:
\begin{itemize}
  \item \texttt{HenselianLocalRing A} does not imply
    \texttt{IsNoetherianRing A}. There are non-Noetherian
    henselian (in particular, henselian valuation) local rings, so
    Noetherianness must be supplied as a hypothesis whenever the
    proof uses Noetherian descent.
  \item \texttt{Algebra.Etale A B} yields
    \texttt{Algebra.FinitePresentation A B} (the standard étale $\Rightarrow$
    finitely presented as an $A$-algebra), but not
    \texttt{Module.FinitePresentation A B}: bridging from finite
    presentation as an $A$-algebra to finite presentation as an
    $A$-module via Mathlib's
    \texttt{Module.FinitePresentation.of\_finite\_of\_finitePresentation}
    would itself consume \texttt{Module.Finite A B}, producing a
    circular instance loop. Hence \texttt{Module.Finite A B} must be
    bundled separately.
\end{itemize}
These two binders match the instance pattern carried by the per-index
single-idempotent lift
\texttt{Algebra.Etale.lift\_idempotent\_henselianPair} (file
\texttt{Proetale/Mathlib/RingTheory/Etale/HenselianPairLift.lean}),
which propagates them through the Newton-iteration carrier. They are
exactly the standard setting in which one decomposes a finite algebra
over a henselian local ring as a finite product of (henselian) local
algebras — the framework of Stacks tags
\href{https://stacks.math.columbia.edu/tag/04GG}{04GG},
\href{https://stacks.math.columbia.edu/tag/04GH}{04GH}, and the
Hensel-idempotent-lift fragment
\href{https://stacks.math.columbia.edu/tag/0DXB}{0DXB}; in the
\href{https://stacks.math.columbia.edu/tag/04GE}{04GE} formulation
the structural theorem states that a finite algebra over a henselian
local ring decomposes as a finite product of finite local henselian
algebras. The signature thus instantiates the canonical Stacks setting
verbatim.

\begin{proof}
  \uses{lem:henselianPair-l3a-idempotent-lift,
        lem:henselianPair-newton-step,
        lem:henselianPair-l1-separation,
        lem:henselianPair-jac}
  \leanok
  Write $\bar B := B / \mathfrak{m} \cdot B$ and let
  $\pi : B \to \bar B$ be the canonical quotient. Composing with the
  base-change isomorphism $\bar B \simeq k \otimes_A B$ identifies
  $\bar B$ with the product $\prod_{i \in I} k_i$ via the algebra
  isomorphism $\mathrm{eqv} : \bar B \xrightarrow{\sim}
  \prod_{i \in I} k_i$. For each $i \in I$ set
  $\bar e_i := \mathrm{eqv}^{-1}(\mathrm{Pi.single}\,i\,1) \in \bar B$,
  the orthogonal idempotents of the product transported back to
  $\bar B$. Each $\bar e_i$ is idempotent because
  $\mathrm{Pi.single}\,i\,1$ is, and $\mathrm{eqv}^{-1}$ is a ring map.
  We also record that
  $\bar e_i \cdot \bar e_j = 0$ for $i \ne j$ and
  $\sum_{i \in I} \bar e_i = 1$, both holding in $\bar B$ because they
  hold componentwise in $\prod_i k_i$.

  \medskip

  \emph{Per-index single-idempotent lift.} By
  \cref{lem:henselianPair-jac} the pair
  $(B, \mathfrak{m} \cdot B)$ is henselian whenever $A$ is henselian
  local and $B$ is étale over $A$, since
  $\mathfrak{m} \cdot B$ sits inside the Jacobson radical of $B$ and
  the underlying root-lifting hypothesis transports along the étale
  structure. Apply
  \cref{lem:henselianPair-l3a-idempotent-lift}
  (Lean carrier \texttt{Algebra.Etale.lift\_idempotent\_henselianPair})
  to each $\bar e_i$ to obtain
  $\tilde e_i \in B$ with $\tilde e_i^2 = \tilde e_i$ and
  $\pi(\tilde e_i) = \bar e_i$.

  \medskip

  \emph{Uniqueness of the Hensel-lift for $Y^2 - Y$.} The argument
  for orthogonality and completeness rests on the following
  uniqueness statement: if $e, e' \in B$ are two idempotents with
  $\pi(e) = \pi(e')$, then $e = e'$. Indeed, the Newton iteration
  for $f(Y) = Y^2 - Y$ starting from any idempotent is stationary
  (an idempotent is already a root of $f$), so both $e$ and $e'$
  arise as the Hensel-pair limit of the same residual data;
  by $\mathfrak{m} \cdot B$-adic separation
  (\cref{lem:henselianPair-l1-separation},
  $\bigcap_n (\mathfrak{m} \cdot B)^n = 0$ in the henselian pair),
  combined with the contractive property of the Newton step
  (\cref{lem:henselianPair-newton-step}), the two limits coincide.
  In particular, an idempotent of $B$ whose residue in $\bar B$
  equals $0$ (respectively $1$) must equal $0$ (respectively $1$).

  \medskip

  \emph{Orthogonality.} Fix $i \ne j$. The element
  $\tilde e_i \tilde e_j$ is an idempotent of $B$
  (product of two commuting idempotents), and
  \[
    \pi(\tilde e_i \tilde e_j)
    = \pi(\tilde e_i) \cdot \pi(\tilde e_j)
    = \bar e_i \cdot \bar e_j
    = 0,
  \]
  the last equality being the componentwise computation in
  $\prod_k k_k$: at index $i$, $1 \cdot 0 = 0$; at index $j$,
  $0 \cdot 1 = 0$; elsewhere, $0 \cdot 0 = 0$. By the uniqueness
  statement above, applied to the idempotents $\tilde e_i \tilde e_j$
  and $0$, we conclude $\tilde e_i \tilde e_j = 0$.

  \medskip

  \emph{Completeness.} The pairwise-orthogonality just established
  makes $\sum_{i \in I} \tilde e_i$ an idempotent of $B$
  (sum of pairwise-orthogonal idempotents). Its residue is
  \[
    \pi\bigl(\textstyle\sum_{i \in I} \tilde e_i\bigr)
    = \textstyle\sum_{i \in I} \pi(\tilde e_i)
    = \textstyle\sum_{i \in I} \bar e_i
    = \mathrm{eqv}^{-1}\bigl(\textstyle\sum_{i \in I}
        \mathrm{Pi.single}\,i\,1\bigr)
    = \mathrm{eqv}^{-1}(1)
    = 1.
  \]
  By the same uniqueness statement applied to
  $\sum_i \tilde e_i$ and $1$, we conclude
  $\sum_{i \in I} \tilde e_i = 1$.

  \medskip

  \emph{Conclusion.} The family $\{\tilde e_i\}_{i \in I}$ is a
  complete orthogonal idempotent system in $B$, and by construction
  each $\tilde e_i$ projects to $\bar e_i$, hence to
  $\mathrm{Pi.single}\,i\,1$ under the residue isomorphism
  $\mathrm{eqv}$. This is the asserted lift.
\end{proof}

\medskip\noindent\textbf{Status.} The theorem is fully formalized.
Iteration~128 closed the six-step backbone of the proof:
(1)~the residue isomorphism $B / \mathfrak{m} B \xrightarrow{\sim}
k \otimes_A B$ via composing
\texttt{Algebra.TensorProduct.quotIdealMapEquivTensorQuot} with
\texttt{Algebra.TensorProduct.comm};
(2)~residue idempotents $\bar e_i := \mathrm{eqv}^{-1}(\mathrm{Pi.single}\,i\,1)$
established as idempotents in $B / \mathfrak{m} B$;
(3)~per-index single-idempotent lifting via
\texttt{Algebra.Etale.lift\_idempotent\_henselianPair} from
\texttt{Proetale/Mathlib/RingTheory/Etale/HenselianPairLift.lean};
(4)~pairwise orthogonality $\tilde e_i \tilde e_j = 0$ for $i \ne j$;
(5)~completeness $\sum_i \tilde e_i = 1$; and
(6)~packaging into \texttt{CompleteOrthogonalIdempotents}. Iteration~129
amends the Lean signature with the two instance binders
\texttt{[IsNoetherianRing A]} and \texttt{[Module.Finite A B]}
documented under \emph{Prerequisites} above; this discharges the two
typed sub-sorries surfaced at the end of iter-128 and brings the file
to a sorry-free state.

\medskip\noindent\textbf{Route note (Newton iteration $\to$ Jacobson
idempotent).} The proof sketch in this chapter motivates orthogonality
and completeness via a uniqueness statement for the Hensel lift of
$Y^2 - Y$, traditionally proved by identifying the Newton-Cauchy
limit (à la Hiblot 1975 / Sharp--Wadsworth 1976) and invoking
$\mathfrak{m} \cdot B$-adic separation. In the Lean formalization
this is replaced by a strictly cleaner Jacobson-idempotent argument:
the auxiliary lemma
\texttt{\_isIdempotentElem\_eq\_zero\_of\_mem\_jacobson\_bot}
(file \texttt{HenselianIdempotentLift.lean}) shows that any idempotent
$e \in B$ lying in the Jacobson radical $\mathrm{jac}(B)$ vanishes,
since then $1 - e$ is a unit, $e \cdot (1 - e) = e - e^2 = 0$, and
cancellation by $(1 - e)$ forces $e = 0$. Applying this to the
candidate idempotents $\tilde e_i \tilde e_j$ (for $i \ne j$) and
$1 - \sum_i \tilde e_i$, both of which are idempotents lying in
$\mathfrak{m} \cdot B \subseteq \mathrm{jac}(B)$, yields orthogonality
and completeness without invoking the Newton-Cauchy limit
identification or the $\mathfrak{m} \cdot B$-adic separation step.
This is a route detail, not a content change: the theorem statement
is unchanged. The Jacobson-idempotent helper is a Mathlib-PR-quality
single-lemma extraction.
